\section{Spatial models and shared transport technology}
\label{ch:equilibrium}

The policy derivative requires two components. A transport block maps
edge--mode costs into routes, modal traffic, and congestion. A spatial model
maps the resulting bilateral costs into economic activity and welfare. The
package shares the transport block across applications and supplies separate
economic-geography and urban equilibrium systems.

\subsection{Two spatial models, one transport block}

The two models start from the same directed network and the same primitive
edge--mode policy. They use the negative-power modal logsum, recursive route
aggregation, flexible and fixed route or mode comparisons, edge and terminal
congestion, and primitive-cost pass-through. They differ in the economic
flows assigned to the network and in the equilibrium response projected onto
welfare.

\begin{longtable}{L{0.20\textwidth}L{0.35\textwidth}L{0.35\textwidth}}
\toprule
Model element & Economic geography & Urban commuting \\
\midrule
\endhead
Economic movement & Goods shipped between locations & Commuters traveling
between residences and workplaces \\
Location margins & Labor and income; source and destination income margins
coincide & Residence \(l^R\) and workplace \(l^F\) margins can differ \\
Network flow & Value traffic normalized by world income & Commuter traffic
normalized by total commuters \\
Spatial response & Wages and the location of mobile labor & Residence and
workplace choices \\
Welfare projection & Common real income or utility under mobility &
\(-d\log\varpi_U/\theta_U\) \\
Route curvature & \(\sigma-1\) & Commuting heterogeneity \(\theta_U\) \\
Typical modes & Road, rail, and water freight & Road, bus, rail or streetcar,
and ferry \\
\bottomrule
\end{longtable}

The table separates economic interpretation from transport technology.
Road, bus, rail, streetcar, and ferry enter as modes in the urban transport
network; public transit is not a third spatial model. Adding a mode changes the
costs and traffic available to commuters without changing the
residence--workplace equilibrium. Conversely, using the same road network in
the two models does not make their welfare derivatives identical because their
spatial states and welfare rows differ.

\subsection{Economic-geography model}

The economic-geography closure turns bilateral transport costs into trade,
production, consumption, and mobile labor. We first define those location-level
objects, then derive bilateral trade and the transformed equilibrium used by
the welfare derivative.

\subsubsection{Production, consumption, and mobility}

The economic-geography model in \citet{AllenFuchsWong2026} has \(N\)
locations. Location \(i\) has labor \(L_i\), wage \(w_i\), productivity
\(A_i\), and amenity \(u_i\). A differentiated variety is produced with labor
under constant returns and perfect competition. Consumers have CES preferences
with elasticity \(\sigma>1\). With \(P_i\) denoting the CES price index,
indirect utility in location \(i\) is proportional to
\[
  U_i=\frac{w_i u_i}{P_i}.
\]

Local activity can affect productivity and amenities:
\[
  A_i=\bar A_i L_i^\alpha,
  \qquad
  u_i=\bar u_i L_i^\beta.
\]
The parameter \(\alpha\) is the elasticity of productivity with respect to
local labor, while \(\beta\) is the corresponding amenity elasticity. Labor is
perfectly mobile across
locations, aggregate labor is fixed, and each worker supplies one unit
inelastically. Spatial equilibrium equalizes utility across occupied
locations and clears goods and labor markets.

These assumptions determine the economic-geography welfare formula. Perfect
mobility equalizes utility, while fixed aggregate labor and labor-only
production exclude further adjustment through factor income or labor supply.
The urban model below replaces these assumptions with separate residence and
workplace choices.

\subsubsection{Bilateral trade and market access}

Let \(\tau_{ij}\) be the composite iceberg cost of delivering a good from
\(i\) to \(j\). CES demand and the production assumptions imply
\[
 X_{ij}
 =W^{1-\sigma}
 \left(\frac{A_i u_j}{\tau_{ij}}\right)^{\sigma-1}
 w_i^{1-\sigma}w_j^\sigma L_j.
\]
Here \(W\) is aggregate welfare in the normalization used by the paper. The
network determines \(\tau_{ij}\); the economic-geography equilibrium
determines wages, labor, productivity, and amenities.

Direct computation of the bilateral cost matrix would obscure the network
structure. Following \citet{AllenFuchsWong2026}, let \(\Pi_i\) denote the
outward price-index analogue and define the inward and outward market-access
stocks as
\[
  \mathcal P_i=P_i^{1-\sigma},
  \qquad
  \mathcal Q_i=\Pi_i^{1-\sigma}.
\]
Define the local terms
\[
 a_i^Q=W^{1-\sigma}u_i^{\sigma-1}w_i^\sigma L_i,
 \qquad
 a_i^P=A_i^{\sigma-1}w_i^{1-\sigma}.
\]
The two market-access stocks satisfy oppositely oriented local recursions:
\begin{align}
 \mathcal Q_i
 &=a_i^Q+\sum_{j\in\mathcal N^+(i)}
       \kappa_{ij}^{1-\sigma}\mathcal Q_j,
 \label{eq:guide-outward-access}\\
 \mathcal P_i
 &=a_i^P+\sum_{j\in\mathcal N^-(i)}
       \kappa_{ji}^{1-\sigma}\mathcal P_j.
 \label{eq:guide-inward-access}
\end{align}
Thus current-node access is divided between local absorption and continuation
through adjacent edges. These sparse recursions replace the dense bilateral
cost matrix. \citet[Appendix A]{AllenFuchsWong2026} derives them from goods
market clearing; Appendix \ref{app:crosswalk} of this guide shows how the
terms in these recursions enter the code.

\subsubsection{Trade exposure and the transformed equilibrium}

Define economic-geography trade exposure by
\[
  \mathcal T_i^G=\frac{\mathcal P_i\mathcal Q_i}{Y^W},
  \qquad
  Y^W=\sum_i w_iL_i.
\]
The access recursions divide this stock into local absorption and continuation.
Define
\[
 s_i^x=\frac{a_i^Q}{\mathcal Q_i},\qquad
 \mu_{ij}=\frac{\kappa_{ij}^{1-\sigma}\mathcal Q_j}{\mathcal Q_i},
 \qquad
 s_i^y=\frac{a_i^P}{\mathcal P_i},\qquad
 \lambda^{in}_{ij}
 =\frac{\kappa_{ij}^{1-\sigma}\mathcal P_i}{\mathcal P_j}.
\]
Here \(\mu_{ij}\) is the outward continuation share on \(i\to j\), while
\(\lambda^{in}_{ij}\) is the same edge's share of destination-\(j\) exposure.
Directed edge traffic factors as
\[
  \Xi_{ij}=\mu_{ij}\mathcal T_i^G
          =\lambda^{in}_{ij}\mathcal T_j^G
          =\frac{\kappa_{ij}^{1-\sigma}
                  \mathcal P_i\mathcal Q_j}{Y^W}.
\]
Outgoing and incoming traffic totals therefore satisfy
\[
 O_i=(1-s_i^x)\mathcal T_i^G,
 \qquad
 I_i=(1-s_i^y)\mathcal T_i^G,
\]
for local absorption shares \(s_i^x\) and \(s_i^y\). The package checks
that outgoing and incoming recovery of \(\mathcal T_i^G\) agree.

\citet[Appendix A]{AllenFuchsWong2026} work with the transformed variables
\[
 x_i=w_i^\sigma L_i^{1+\beta(\sigma-1)},\qquad
 y_i=W^{\sigma-1}w_i^{1-\sigma}L_i^{\alpha(\sigma-1)}.
\]
This transformation separates sparse market-access propagation from the
aggregate labor constraint. Define
\[
 e=1+\beta(\sigma-1)+\alpha\sigma,
 \qquad
 \rho=\frac{1+\alpha+\beta}{e}.
\]
The inverse map from \((x,y)\) to wages, labor, and welfare requires
\(e\neq0\); the welfare normalization also requires
\(1+\alpha+\beta\neq0\).
Stack the transformed log changes in the \(2N\)-element column vector
\[
 z_G=(d\log x_1,\ldots,d\log x_N,
      d\log y_1,\ldots,d\log y_N)^\top.
\]
The route, mode, and congestion maps below translate a primitive edge--mode
shock into the aggregate edge-cost changes entering the residual system for
\(z_G\).

\implementationbox{The economic-geography coefficient map is implemented by
\code{AdjointRSUE.coefs}. \code{AdjointRSUE.assemble\_J} constructs the
no-congestion spatial Jacobian. The generic builder then adds the common
route, modal, and congestion feedback.}

\subsection{Urban commuting: residences and workplaces}

The urban specification follows the residence--workplace model in
\citet{AA_2022restud}. Let \(l_i^R\) and \(l_i^F\) denote normalized residence
and workplace masses. A location can contain many residents but few jobs, or
many jobs but few residents. Commuting flows therefore have distinct origin
and destination margins.

The urban accounting stock satisfies
\[
 \mathcal T_i^U
 =l_i^F+\sum_j\Xi_{ij}
 =l_i^R+\sum_j\Xi_{ji}.
\]
The first equality combines workplace absorption with outgoing route traffic;
the second combines residence absorption with incoming traffic. It is the
urban counterpart to trade exposure, but it is measured in normalized
commuter flows rather than shares of world income. The loader rejects a
baseline when the two expressions disagree beyond the declared tolerance.

The differentiated urban state and welfare change are
\[
 z_U=(d\log l^R,d\log l^F,d\log\varpi_U)^\top,
 \qquad
 d\log W=-\frac{d\log\varpi_U}{\theta_U}.
\]
We write the urban equilibrium normalizer as \(\varpi_U\) to distinguish it
from the primitive-cost pass-through \(\chi_{klm}\). The implementation retains
\code{chi\_U} as the historical name of this final state component.
The parameter \(\theta_U>0\) governs commuting heterogeneity and supplies the
route curvature. The urban parameters \(\alpha\) and \(\beta\) enter the
workplace- and residence-side response coefficients. The resulting Jacobian
\(J_U^0\), aggregate-cost loading \(B_{\kappa,U}\), and welfare row
\[
 q_U^\top=(0,\ldots,0,-1/\theta_U)
\]
are specific to the residence--workplace equilibrium. They cannot be obtained
by relabeling the trade equations.

Public transit enters through the shared edge--mode network. A directed edge
can carry road, bus, rail, streetcar, or ferry traffic, and the modal logsum
determines how their generalized costs combine. The urban model then maps
the induced commuting-cost changes into residence, workplace, and welfare
responses. Subsection \ref{ch:extensions} documents the synthetic urban test and
the Seattle data candidate.

\subsection{Interface with the transport block}

For spatial model \(c\in\{G,U\}\), write its no-congestion residual as
\[
 G_c^0(z_c,\kappa)=0.
\]
Each model supplies three inputs to the shared computation: a spatial
Jacobian \(J_{c,0}\), a loading \(B_{\kappa,c}\) from aggregate edge costs into
the residuals, and a welfare row \(q_c\). The next subsection constructs the
route, modal, and congestion derivatives combined with these inputs.
Subsection \ref{ch:ift} then derives welfare after eliminating the transport
block.
